\section{dcbench Framework}

\texttt{dcbench} is a benchmark toolkit designed for data-centric AI tasks. It provides a structured way to define, solve, and evaluate tasks where the focus is on data quality and discovery rather than just model training.

The Slice Discovery task in \texttt{dcbench} aims to identify coherent subgroups of data (slices) where a model performs poorly or behaves unexpectedly.

\begin{table}[h]
\centering
\begin{tabular}{|l|l|l|}
\hline
\textbf{Category} & \textbf{Data Set} & \textbf{Characteristics} \\ \hline
\textbf{Correlation} & CelebA & Performance drops driven by wrong correlations of attributes (gender, age). \\ \hline
\textbf{Rareness} & ImageNet & Low accuracy on infrequent classes. \\ \hline
\textbf{Label Noise} & ImageNet & Apparent errors caused by mislabeling or class ambiguity. \\ \hline
\end{tabular}
\end{table}

\subsection{Problem Formulation}
Let $\mathcal{D} = \{x_i, y_i, \hat{y}_i, z_i\}_{i=1}^n$ be a dataset where:
\begin{itemize}
    \item $x_i$ is the input (e.g., an image).
    \item $y_i$ is the ground truth label.
    \item $\hat{y}_i = f(x_i)$ is the model prediction.
    \item $z_i$ is a latent representation or embedding (e.g., CLIP embeddings or model activations).
\end{itemize}

The dataset also contains hidden ground truth slices $\mathcal{S} = \{S_1, S_2, \dots, S_k\}$, where each $S_j \in \{0, 1\}^n$ is a binary indicator vector marking membership in the $j$-th slice.

A Slice Discovery algorithm (Solver) $\mathcal{A}$ must produce a set of predicted slices:
\[ \hat{\mathcal{S}} = \mathcal{A}(\mathcal{D}) = \{\hat{S}_1, \hat{S}_2, \dots, \hat{S}_m\} \]
where $\hat{S}_j \in [0, 1]^n$ represents the probability or confidence of each instance belonging to the $j$-th discovered slice.

\subsection{Case Studies}

In \texttt{dcbench}, slice discovery problems are categorized based on the nature of the model's failure. We illustrate the setup using three representative case studies.

\subsubsection{General Setup for Case Studies}
For each case study, we define:
\begin{itemize}
    \item \textbf{Classifier Task}: The primary objective of the model $f$, which produces predictions $\hat{y}_i$ for a target $y_i$.
    \item \textbf{Ground Truth Slice ($S_j$)}: A known subgroup where the model's accuracy $\mathbb{P}(\hat{y}_i = y_i | i \in S_j)$ is significantly lower than its overall accuracy.
    \begin{itemize}
        \item \textit{Example 1 (Correlation Slices)}: CelebA (\texttt{p\_72776}), Target: \texttt{wearing\_lipstick}.
        \begin{itemize}
            \item \textbf{Overall Accuracy}: 90.68\%
            \item \textbf{Slice $S_0$}: \texttt{wearing\_lipstick=0\_young=1} (Young people not wearing lipstick).
            \item \textbf{Slice $S_0$ Size}: 19.37\% (1205 / 6221)
            \item \textbf{Slice $S_0$ Accuracy}: 75.93\%
        \end{itemize}
        \item \textit{Example 2 (Rareness Slices)}: ImageNet (\texttt{p\_118818}), Target: \texttt{food.n.01}.
        \begin{itemize}
            \item \textbf{Overall Accuracy}: 82.38\%
            \item \textbf{Slice $S_0$}: \texttt{menu.n.02} (Menus appearing in food contexts).
            \item \textbf{Slice $S_0$ Size}: 1.70\% (103 / 6057)
            \item \textbf{Slice $S_0$ Accuracy}: 68.93\%
        \end{itemize}
        \item \textit{Example 3 (Label Noise Slices)}: ImageNet (\texttt{p\_122446}), Target: \texttt{cat.n.01}.
        \begin{itemize}
            \item \textbf{Overall Accuracy}: 83.36\%
            \item \textbf{Slice $S_0$}: \texttt{cougar.n.01} (Cougars often confused with house cats).
            \item \textbf{Slice $S_0$ Size}: 4.38\% (263 / 6004)
            \item \textbf{Slice $S_0$ Accuracy}: 57.03\%
        \end{itemize}
    \end{itemize}
    \item \textbf{Slice Discovery Goal}: The objective of the solver $\mathcal{A}$ to identify a predicted slice $\hat{S}_l$ that maximizes the similarity to $S_j$.
\end{itemize}


\subsection{Evaluation Metrics}

The quality of the discovered slices is evaluated by comparing them to the ground truth slices. 
For each ground truth slice $S_j \in \mathcal{S}$, we identify the best-matching predicted slice 
$\hat{S}_l \in \hat{\mathcal{S}}$ and compute several metrics.

\noindent\textbf{Precision@k}

\[
P@k(S_j, \hat{S}_l) 
= \frac{1}{k} 
\sum_{i \in \text{top-k}(\hat{S}_l)} \mathbb{I}(i \in S_j)
\]

Common values are $k = 10$ and $k = 25$.

\noindent\textbf{Mean Problem Score}

The final score for a problem averages the best value for each ground truth slice:

\[
\text{Score} = \frac{1}{k} 
\sum_{j=1}^{k} 
\max_{\hat{S}_l \in \hat{\mathcal{S}}}
\text{Metric}(S_j, \hat{S}_l)
\]

\noindent\textbf{Average Precision (AP)}

In the implementation, AP is computed using the slice probability vector 
$\texttt{slice\_prob}$, not binary predictions.

\begin{verbatim}
"average_precision": skmetrics.average_precision_score(
    y_true=slc,
    y_score=slice_prob
)
\end{verbatim}

Given a ground-truth binary indicator $S \in \{0, 1\}^N$ and predicted probabilities $\hat{S} \in [0, 1]^N$, we first sort the
  $N$ samples by their predicted scores in descending order:
  \[ \hat{S}{(1)} \ge \hat{S}{(2)} \ge \dots \ge \hat{S}{(n)} \ge \dots \ge \hat{S}{(N)} \]


  The AP is then defined as:
  \[ AP = \sum_{n=1}^{N} (R_n - R_{n-1}) P_n \]


  where:
  \begin{itemize}
      \item $n$ is the rank index of the samples after sorting by probability.
      \item $P_n$ and $R_n$ are the precision and recall calculated by setting the classification threshold to the $n$-th highest probability score, $T = \hat{S}_{(n)}$.
      \item By convention, $R_0 = 0$.
  \end{itemize}


\subsubsection{Case Study: Correlation Slices}
\noindent\textbf{Example (Correlation Slices)}: CelebA (\texttt{p\_72776}), Target: \texttt{wearing\_lipstick}.
\begin{itemize}
    \item \textbf{Overall Accuracy}: 90.68\%
    \item \textbf{Slice $S_0$}: \texttt{wearing\_lipstick=0\_young=1} (Young people not wearing lipstick).
    \item \textbf{Slice $S_1$}: \texttt{wearing\_lipstick=1\_young=0} (Non-young people wearing lipstick).
    \item \textbf{Slice $S_0$ Size}: 19.37\% (1205 / 6221)
    \item \textbf{Slice $S_1$ Size}: 18.76\% (1167 / 6221)
    \item \textbf{Slice $S_0$ Accuracy}: 75.93\%
    \item \textbf{Slice $S_1$ Accuracy}: 98.11\%
\end{itemize}

In \texttt{dcbench}, the matching between your predicted slices and the ground-truth hidden slices is an all-to-all exhaustive search and the matched slice is picked by maximum average precision.

\noindent\textbf{Evaluation Results}

\begin{table}[h!]
\centering
\begin{tabular}{lcc}
\hline
\textbf{Metric} & \textbf{Slice $S_0$} & \textbf{Slice $S_1$} \\
\hline
pred\_slice\_idx    & 0                        & 5 \\
average precision  & 0.46             & 0.37 \\
precision-at-10     & 0.6                      & 0.4 \\
precision-at-25     & 0.56                     & 0.24 \\
recall              & 0.44                  & 1.00 \\
\hline
\end{tabular}
\end{table}

\begin{figure}[h!]
\centering
\includegraphics[width=0.4\textwidth]{figures/discovered_slice_0.png}
\caption{Discovered Slice 0 (Size: 1260) for CelebA (\texttt{p\_72776}). Example instance (ID: 143706.jpg) has target 0, predicted probabilities [0.1603, 0.8397], and slice membership [True, False].}
\label{fig:discovered_slice_0}
\end{figure}

\noindent\textbf{Average Precision (AP)}

In the implementation, AP is computed using the slice probability vector 
$\texttt{slice\_prob}$, not binary predictions.

\begin{verbatim}
"average_precision": skmetrics.average_precision_score(
    y_true=slc,
    y_score=slice_prob
)
\end{verbatim}

Given a ground-truth binary indicator $S \in \{0, 1\}^N$ and predicted probabilities $\hat{S} \in [0, 1]^N$, we first sort the
  $N$ samples by their predicted scores in descending order:
  \[ \hat{S}{(1)} \ge \hat{S}{(2)} \ge \dots \ge \hat{S}{(n)} \ge \dots \ge \hat{S}{(N)} \]


  The AP is then defined as:
  \[ AP = \sum_{n=1}^{N} (R_n - R_{n-1}) P_n \]


  where:
  \begin{itemize}
      \item $n$ is the rank index of the samples after sorting by probability.
      \item $P_n$ and $R_n$ are the precision and recall calculated by setting the classification threshold to the $n$-th highest probability score, $T = \hat{S}_{(n)}$.
      \item By convention, $R_0 = 0$.
  \end{itemize}


\subsection{Observations}
\begin{itemize}
    \item Their problem formulation is flexible.
    \item $m$ is not necessarily equal to $k$
\end{itemize}